\section{Results}
\subsection{Environmental Impact}


\subsubsection{GHG Emissions assessment}

\paragraph{Regional Energy Mix}

\paragraph{Local Energy Mix}

\begin{figure}[h!]
    \centering
    \includegraphics[width=0.8\textwidth]{figures/placeholder.png}
    \caption{Local energy mix evolution over time.}
    \label{fig:plant_schematic}
\end{figure}

\paragraph{Residual Fossil Fuel Consumption}

\paragraph{Transportation Sector}

mettere commento su miglioramento della qualità dell'aria

\paragraph{Overall GHG Emissions}

\begin{figure}[h!]
    \centering
    \begin{minipage}{0.48\textwidth}
        \centering
        \includegraphics[width=\textwidth]{figures/ghg_emissions.png}
        \caption{GHG emissions over time.}
        \label{fig:ghg_emissions}
    \end{minipage}\hfill
    \begin{minipage}{0.48\textwidth}
        \centering
        \includegraphics[width=\textwidth]{figures/net_electricity_demand.png}
        \caption{Local Net Electricity and Fossil Demand evolution in time.}
        \label{fig:net_electricity_demand}
    \end{minipage}
\end{figure}

\subsubsection{H_$2$O Usage}
The water used for the condensers and the hydrogen production are surely negligible with respect to the mass of water hosted in the neighboring Lago Maggiore, so it's possible to say that the hydric resources of the area are not compromised by the realization of the hybrid energy system proposed.

\subsubsection{Land Use}
The well-established presence of nuclear facilities in the Ispra JRC area allows to state the the land use of this projects is virtually non existent.

\subsection{Conclusions}
The project states that hybrid energy systems can be a viable way to reduce by a great amount the carbon emissions of the investigated area, eliminating all the 
fossil consumption for elecricity demand and paving the way to the reduction of the transportation carbon footprint too. The dymola model of the plant proves the technical feasibilty of the plant and gives an accurate hourly account of all the energetic and material quantities exchanged both inside the boundaries of the plant and through them towards the community, allowing an exact evaluation of the performances of the facilty for economical analysis. The plant layout was designed and dimensioned to let the gas turbine operate at no less that 60 percent of its nominal power to avoid pollutant emissions and inefficiencies: the number of HTSE and SMRs was, in fact, chosen accordingly within the constraints of TANDEM \cite{Tandem} components' sizes. This aspect led to an overdimensioning of the power plant in relation to the zonal elecricity demand, meaning that a great portion of the electricity produced by the nuclear plant is dedicated to noble decarbonizing processes (i.e. the hydrogen and methane production) way less profitable than direct electrical energy sell to the grid.  One of the central outcomes of the economical investigation is that huge capital costs of the system could not be recoverable if the whole hybrid plant were to be constructed immediately, but they could be balanced and overrun by revenues if a progressive implementation of its different sections is pursued. To pay off these high initial costs and wait some years for taking benefit from the further techological maturity of the systems employed, the delay of the installation of the HTSE-Methanator-GasTurbine complex is delayed to the 20th year of operation. During this delay the entire excess electricity produced by the SMRs is directly sold and more profitably sold to the market in order to repay a portion of the nuclear capital costs. This strategy is realistic - resemblant of the way a multiple SMR plant installs one reactor at a time to finance the next one - and able to largely justify the feasibility of the project. The proven economic feasibility conjugated to the strong decarbonizing impact and the the cease of foreing energy import make the project a promising idea for the Sette Laghi region and a wise use of the ESSOR nuclear facility.  

\section{Possible Improvements}
\subsection{On the Data Collection}
It is defenitly possible and could give better results if the data for each class was obtained by a set of real building from the region instead of just one.
Ideally one would use a statistically signfinicant set of buldings and vary the utilization profile as well.
Moreover, it would be beneficial to wheight the simulation data on the real energy requirments of those same buildings,
which is a much more complex task since it would require extensive data acquisition during the year. 
While this is feasible for electrical demand (most operators give to the customer a detailed bill with hourly consumption),
it is much more difficult to get the same data for fossil fuel consumption.
At last, the normalization has been done with the heated area an not with the volume (which is much more indicative of the energy requirment of a building) 
given the on-field knowledge that the volume is usually approximatly obtained by a the rule of thumb of multiplying the heated area by a factor of 2.7 or 3 in most cases.
From the simulation we have neglected humidification and dehumidification electricity needs since these plants are rearly installed today,
a better evaluation would evaluate the possible future share of air treatment plants in the residential sector.

Photovoltaic: We have neglected any solar-thermal system installed since their are considered not very impactful and their impact is complex to evaluate.

Convertion thermal to primary: use time dependent efficiency for heat pumps that take into account the ambient temperature,
this could be accomplished considering that our code uses hourly data.

Other Electricity Demand: This is without a doubt the most complex contribution to evaluate,
expecially considering that it is not easy to imagine how the electricity consumption will evolve in 25 years.
It could be improved by collecting a statistically signficant number of examples from the region on the type of appliances installed and their consumption.

Overall: the methodology of using class distribution as a way of evaluating the energy demand is considered a good approach,
for the best results we shall gather consumption data for both electric and fossil fuels for a statistically significant number of buildings for each energy class.

\subsection{On the Plant Modelling}
\subsection{On the Economic Study}

\section{Conclusion}
